% part: intuitionistic-logic
% chapter: semantics
% section: topological-semantics

\documentclass[../../../include/open-logic-section]{subfiles}

\begin{document}

\olfileid{int}{sem}{top}

\olsection{સ્થલરચનાત્મક અર્થવિજ્ઞાન}

સ્ફુરણાવાદી તર્ક માટે અર્થવિજ્ઞાન આપવાની બીજી રીત
સ્થલરચનાના ગાણિતિક ખ્યાલનો ઉપયોગ કરવાની છે.

\begin{defn}
  ધારો કે $X$ એક ગણ છે. $X$ પરની \emph{સ્થલરચના} એ
  $\Top{O} \subseteq \Pow{X}$ એવો ગણ છે જે નીચેના ગુણધર્મો
  સંતોષે છે. $\Top{O}$ના !!{element}sને સ્થલરચનાના
  \emph{ખુલ્લા ગણો} કહે છે. $X$ અને $\Top{O}$ને સાથે મળીને
  \emph{સ્થલરચનાત્મક અવકાશ} કહે છે.
  \begin{enumerate}
  \item ખાલી ગણ અને આખો અવકાશ ખુલ્લા છે: $\emptyset$, $X \in
    \Top{O}$.
  \item ખુલ્લા ગણો સાન્ત છેદ હેઠળ બંધ છે: જો $U$, $V \in
    \Top{O}$ હોય, તો $U \cap V \in \Top{O}$.
  \item ખુલ્લા ગણો મનસ્વી યોગગણ હેઠળ બંધ છે: જો બધા $i \in I$
    માટે $U_i \in \Top{O}$ હોય, તો $\bigcup \Setabs{U_i}{i \in
      I} \in \Top{O}$.
  \end{enumerate}
\end{defn}

ખુલ્લા ગણોનો સંગ્રહ સંદર્ભમાંથી જાણી શકાય ત્યારે આપણે
સ્થલરચનાત્મક અવકાશ માટે માત્ર $X$ લખી શકીએ છીએ. જોકે,
$X$ પર ખુલ્લા ગણો નક્કી કર્યા પછી જ તેને સ્થલરચનાત્મક
અવકાશ કહી શકાય.
\sourcecorrection{OLMD-087}{મૂળ આ વાક્યમાં $X$ને
``સ્થલરચના'' કહે છે; ઉપરની વ્યાખ્યા મુજબ સ્થલરચના
$\Top{O}$ છે, જ્યારે $X$ અને $\Top{O}$ સાથે મળીને
સ્થલરચનાત્મક અવકાશ બને છે.}

\begin{defn}
  સ્ફુરણાવાદી વિધાનાત્મક તર્કનું \emph{સ્થલરચનાત્મક નિદર્શ}
  એ ત્રિક $\mModel{X} = \tuple{X, \Top{O}, V}$ છે, જ્યાં
  $\Top{O}$ એ~$X$ પરની સ્થલરચના છે અને $V$ એ એવું વિધેય
  છે જે દરેક વિધાનાત્મક ચલને $\Top{O}$નો એક ખુલ્લો ગણ આપે છે.

  આપેલું સ્થલરચનાત્મક નિદર્શ~$\mModel{X}$ હોય, તો
  $\Prop{X}{!A}$ને નીચે મુજબ આગમનાત્મક રીતે વ્યાખ્યાયિત કરીએ છીએ:
  \begin{enumerate}
  \item $\Prop{X}{\lfalse} = \emptyset$
  \item $\Prop{X}{p} = V(p)$
  \item $\Prop{X}{!A \land !B} = \Prop{X}{!A} \cap \Prop{X}{!B}$
  \item $\Prop{X}{!A \lor !B} = \Prop{X}{!A} \cup \Prop{X}{!B}$
  \item $\Prop{X}{!A \lif !B} = \Interior{(X \setminus \Prop{X}{!A}) \cup
    \Prop{X}{!B}}$
  \end{enumerate}
  અહીં $\Interior{V}$ એ $V \subseteq X$ ગણને તેના
  \emph{અંતર્ભાગ} પર લઈ જતું વિધેય છે; એટલે, તેમાં સમાયેલા
  બધા ખુલ્લા ગણોનો યોગગણ. બીજા શબ્દોમાં,
  \[
  \Interior{V} = \bigcup \Setabs{U}{U \subseteq V \text{ અને } U \in \Top{O}}.
  \]
\end{defn}

નોંધો કે કોઈ પણ ગણનો અંતર્ભાગ હંમેશા ખુલ્લો છે, કારણ કે
તે ખુલ્લા ગણોનો યોગગણ છે. તેથી $\Prop{X}{!A}$ હંમેશા
ખુલ્લો ગણ છે.

સ્થલરચનાત્મક અર્થવિજ્ઞાન ખૂબ અમૂર્ત છે, છતાં તેને સમજવાની
કેટલીક રીતો તેની પાછળનો વિચાર સ્પષ્ટ કરી શકે છે. ધારો કે
$X$ના !!{element}s, અથવા ``બિંદુઓ'', એવા બિંદુઓ છે જ્યાં
વિધાનોનું મૂલ્યાંકન કરી શકાય છે. $!A$ જ્યાં સાચું હોય એવા
બધા બિંદુઓનો ગણ એ~$!A$ દ્વારા વ્યક્ત થયેલું વિધાન છે.
બિંદુઓનો દરેક ગણ શક્ય વિધાન નથી; ફક્ત $\Top{O}$ના
!!{element}s જ શક્ય વિધાનો છે. દરેક~$X$ માટે $!A \Entails !B$
ત્યારે અને માત્ર ત્યારે જ જ્યારે $!A$ જ્યાં સાચું હોય તે
દરેક બિંદુએ $!B$ પણ સાચું હોય, એટલે $\Prop{X}{!A}
\subseteq \Prop{X}{!B}$. અસંગત વિધાન~$\lfalse$ ક્યારેય
સાચું નથી, તેથી $\Prop{X}{\lfalse} = \emptyset$.

$!B \land !C$, $!B \lor !C$ અને $!B \lif !C$ દ્વારા
વ્યક્ત થયેલાં વિધાનોનો, $!B$ અને~$!C$ દ્વારા વ્યક્ત થયેલાં
વિધાનો સાથે કયો સંબંધ હોવો જોઈએ જેથી સ્ફુરણાવાદી રીતે
પ્રામાણ્ય નિયમો જળવાય, એટલે $!A \Proves !B$ ત્યારે અને માત્ર
ત્યારે જ જ્યારે $\Prop{X}{!A} \subseteq \Prop{X}{!B}$?
\sourcecorrection{OLMD-088}{મૂળમાં આ સમકક્ષતા અને અંતિમ
ફકરાની શરતમાં $\subset$ છે; સમાન ગણો માટે પણ અનુસરણ
શક્ય હોવાથી બંને જગ્યાએ $\subseteq$ કર્યું છે.}
કોઈ પણ $!A$ માટે $\lfalse \Proves !A$ જોઈએ, અને બધા
$U$ માટે $\emptyset \subseteq U$ હોવાથી આ શરત સંતોષાય
છે. $!B \land !C \Proves !B$ હોવાથી
$\Prop{X}{!B \land !C} \subseteq \Prop{X}{!B}$ જોઈએ;
એ જ રીતે $\Prop{X}{!B \land !C} \subseteq \Prop{X}{!C}$
જોઈએ. $W \subseteq U$ અને $W \subseteq V$ સંતોષતો સૌથી
મોટો ગણ $U \cap V$ છે. ઊલટું, $!B \Proves !B \lor !C$
અને $!C \Proves !B \lor !C$ હોવાથી
$\Prop{X}{!B} \subseteq \Prop{X}{!B \lor !C}$ અને
$\Prop{X}{!C} \subseteq \Prop{X}{!B \lor !C}$ જોઈએ.
$U \subseteq W$ અને $V \subseteq W$ સંતોષતો સૌથી નાનો
ગણ~$W$ એ $U \cup V$ છે.

$\lif$ માટેની વ્યાખ્યા થોડી મુશ્કેલ છે: $!A \lif !B$ એ
એવું સૌથી નબળું વિધાન વ્યક્ત કરે છે જે $!A$ સાથે મળીને
$!B$ને ફલિત કરે. $!A \lif !B$ અને $!A$ સાથે મળીને~$!B$
ફલિત થાય છે તે $(!A \lif !B) \land !A \Proves !B $
પરથી સ્પષ્ટ છે. તેથી $\Prop{X}{!A \lif !B}$ એવો સૌથી મોટો
ખુલ્લો ગણ હોવો જોઈએ કે $\Prop{X}{!A \lif !B} \cap \Prop{X}{!A}
\subseteq \Prop{X}{!B}$; આથી આપણી વ્યાખ્યા મળે છે.

\end{document}
